\chapter{Conclusão}\label{cap:conclusao}

Esta dissertação partiu de uma análise da \emph{baseline} do desafio SoccerNet Game State Reconstruction, que revelou dois pontos fracos na identificação dos jogadores. O primeiro era o reconhecimento do número da camisola, correto em apenas $15.3\%$ dos \emph{tracklets}. O segundo era a atribuição da equipa e do respetivo lado do campo, que falhava em $28$ dos $159$ vídeos avaliados. Foi sobre estes dois problemas que o trabalho se concentrou.

Para os resolver, foram alterados quatro módulos. O módulo de reconhecimento do número da camisola foi substituído, enquanto os módulos de atribuição da equipa e do lado do campo foram reescritos. A reidentificação também foi modificada, mas apenas para calcular um descritor de cor utilizado no agrupamento das equipas, sem alterar o que já produzia. Os restantes cinco módulos foram mantidos idênticos aos da \emph{baseline}. São a deteção, o \emph{tracking}, a estimativa dos pontos do campo, a calibração e a agregação por \emph{tracklet}.

Manter estes cinco módulos inalterados foi uma decisão de método, e não apenas uma delimitação do âmbito do trabalho. Como todos os restantes componentes permaneceram iguais, qualquer diferença de desempenho entre os dois sistemas pode ser atribuída às alterações introduzidas. Com o mesmo objetivo, a \emph{baseline} foi executada nos mesmos $159$ vídeos, no mesmo ambiente e sobre os mesmos \emph{tracklets}. Restava ainda separar o efeito de cada alteração. Para isso, o GS-HOTA foi decomposto por atributo.

É neste ponto que o trabalho se distingue da maioria das soluções apresentadas ao desafio. Estas modificam vários módulos ao mesmo tempo e, por isso, mesmo quando o GS-HOTA sobe muito, é difícil saber quanto do ganho resulta de cada alteração. Nesta dissertação, graças à decomposição por atributo, esse efeito lê-se de forma direta.

No reconhecimento do número da camisola, o OCR da \emph{baseline} foi substituído pelo modelo de visão e linguagem Qwen2.5-VL-3B. Para caber nos 6~GB de memória da placa gráfica, o modelo foi quantizado a 4 bits. A substituição do modelo não foi, contudo, a única alteração. Mudou também a forma como o problema lhe é apresentado, em três passos.

Primeiro, em vez de lhe ser entregue o recorte completo do jogador, isola-se a região do tronco a partir da pose estimada. Depois, descartam-se os recortes sem informação útil. Por fim, o número de cada \emph{tracklet} é decidido apenas com as leituras em que a confiança do modelo, medida pela probabilidade da resposta, passa um limiar.

Este último passo resultou de um diagnóstico da primeira versão do módulo, que tratava todas as leituras como igualmente fiáveis. Com a introdução deste critério, os erros reduziram-se para cerca de um terço, passando de $12.3\%$ para $4.4\%$ das deteções de jogadores, enquanto os acertos aumentaram de $51.4\%$ para $57.0\%$.

Na atribuição da equipa, as alterações abrangeram as duas fases do módulo, o agrupamento dos jogadores em duas equipas e a escolha do lado de cada uma. No agrupamento, quando a reidentificação não consegue distinguir as duas equipas, os jogadores passam a ser agrupados pela cor da camisola. Essa cor é medida na região do tronco, durante a própria reidentificação. O agrupamento também passou a validar o seu resultado. Se uma das equipas ficar com menos de $8\%$ das deteções, o agrupamento é rejeitado e refeito com base na cor.

Na escolha do lado, o critério único da \emph{baseline}, baseado na posição média de cada equipa, deu lugar a uma cascata de quatro critérios, aplicados por ordem até que um deles decida. Cada critério é menos exigente do que o anterior, e o último consiste num desempate pelo guarda-redes.

Os ganhos foram medidos em todos os vídeos anotados, por comparação com a \emph{baseline} executada nas mesmas condições. No número da camisola, a exatidão por \emph{tracklet} subiu de $15.3\%$ para $42.6\%$ nos $159$ vídeos. Houve melhoria em $156$ vídeos e nenhum piorou. Na atribuição da equipa e do lado, a exatidão passou de $83.5\%$ para $92.8\%$ dos \emph{tracklets}. Na inspeção visual, os vídeos em que o módulo de equipa falha desceram de $28$ para $7$.

Na métrica oficial do desafio, o GS-HOTA aumentou $16.57$ pontos no conjunto de teste, de $31.30\%$ para $47.87\%$, $13.11$ na validação e $13.85$ no treino. A decomposição por atributo mostra que todo este ganho resulta dos dois atributos alterados, a equipa e o número. Entre os dois, é o número da camisola que tem maior peso. O ganho é de magnitude semelhante nos três subconjuntos, pelo que não depende de um subconjunto em particular.

Para comparar os dois sistemas, foi ainda necessário corrigir um defeito da \emph{baseline}. O GS-HOTA calculado apenas com base na posição deveria ser igual nos dois sistemas, mas não era. A causa residia na configuração distribuída da \emph{baseline}, que calibra a câmara em lotes de oito \emph{frames} e aplica a homografia do primeiro aos restantes sete. No conjunto de teste, este defeito custava à \emph{baseline} cerca de $5$ pontos de GS-HOTA na posição e $2.4$ pontos na métrica completa. Todos os valores da \emph{baseline} apresentados nesta dissertação correspondem à versão corrigida, de modo a que o ganho reportado resulte apenas das alterações introduzidas. A correção não constitui um contributo deste trabalho, mas pode ser útil a quem utilize a \emph{baseline} na configuração por omissão.

Delimitar o alcance destes ganhos ocupou uma parte importante do trabalho, e começou antes da escolha do modelo. Para a leitura do número, foram comparadas várias abordagens (Tabela~\ref{tab:alternativas}). Entre elas estavam o EasyOCR, o PaddleOCR afinado, o HunyuanOCR, a versão de 7B do Qwen e a super-resolução das imagens. Todas foram postas de parte, porque não liam melhor, não cabiam nos 6~GB da placa gráfica ou eram demasiado lentas.

Já com o Qwen escolhido, foram testadas três formas de melhorar a leitura. A primeira ponderava as leituras pela qualidade do recorte, a segunda realçava o contraste e a terceira juntava todas as vistas de um \emph{tracklet} numa só imagem. Nenhuma melhorou o resultado, e o realce de contraste chegou a piorar vídeos que já estavam certos. As três esbarraram no mesmo obstáculo. Nos casos difíceis, o número não está visível no recorte, pela distância da câmara ou pela orientação do jogador, e nenhuma transformação da imagem o recupera. O que funcionou foi outra via, usar a confiança do modelo para separar as leituras em que viu o número daquelas em que o adivinhou.

Depois disso, foram testadas mais três correções, e todas foram rejeitadas. A fusão de fragmentos do \emph{tracking} subia a exatidão por \emph{tracklet}, mas quase não mudava o GS-HOTA. Recuperar as leituras certas com pouca confiança trazia de volta os números inventados que o limiar tinha cortado. A exigência de que cada número fosse único dentro de cada equipa apagava sobretudo leituras certas, porque a maioria dos conflitos eram fragmentos do mesmo jogador.

Na atribuição da equipa, a inspeção visual dos $159$ vídeos mostra que as falhas que restam têm duas origens. A primeira é o classificador de função, que não foi alterado. Em $20$ vídeos da validação, este classificador dá os jogadores de amarelo como árbitros nos dois sistemas, e essa equipa fica de fora do agrupamento. É isso que deixa a validação abaixo dos outros dois conjuntos. A segunda origem é o próprio módulo de equipa, que ainda erra em sete vídeos, contra $28$ na \emph{baseline}. Em três, o lado está trocado. Nos outros três, há um \emph{tracklet} na equipa errada. No último, as duas equipas ficam num só grupo, porque a verificação do agrupamento não o conseguiu refazer pela cor.

Para perceber até onde a posição dos jogadores indica o lado de cada equipa, analisou-se também a anotação dos $160$ vídeos. Esta análise inclui o SNGS-096, que fica fora da comparação dos sistemas. Em cada vídeo, compararam-se dois critérios com o lado anotado. O primeiro é a posição do guarda-redes, que fica junto à sua baliza e nunca contradiz o lado anotado da sua equipa, nos $156$ vídeos em que aparece. O segundo é a posição média dos jogadores de campo, o critério usado pela generalidade das soluções do desafio. Este falha em sete vídeos, por dois motivos diferentes. Em dois, as duas equipas ocupam quase a mesma zona do campo, e as suas posições médias diferem menos de um metro. Nos outros cinco, o erro está na própria anotação, que dá aos jogadores de campo o lado oposto ao do seu guarda-redes.

Estes cinco vídeos mostram também os limites da métrica de equipa. Neles, um sistema que atribua o lado certo tem $0\%$ de exatidão, porque a anotação diz o contrário. Além disso, a exatidão por \emph{tracklet} não contabiliza os jogadores classificados como árbitros nem os \emph{tracklets} que trocam de jogador, pelo que não capta parte dos erros. Os valores foram mantidos tal como a métrica os conta, para garantir a comparação com outros trabalhos. Ainda assim, uma análise fina desta tarefa exige as duas fontes, a métrica e a inspeção visual.

\section{Trabalho Futuro}\label{sec:trabalho_futuro}

Os resultados apontam para seis linhas de trabalho futuro, apresentadas por ordem de prioridade.

\begin{enumerate}
    \item \textbf{Medição da cor do fundo em cada recorte.} O descritor de cor atual exclui sempre a mesma gama de tons, assumindo que o verde corresponde à relva. Esta suposição falha quando uma das equipas joga com essa cor, como no SNGS-069, de verde e branco. Medir o fundo a partir dos píxeis em redor de cada jogador resolveria este caso. A alteração é local, e os módulos de equipa correm em menos de dois segundos por vídeo.

    \item \textbf{Maior peso do guarda-redes na escolha do lado.} Atualmente, o guarda-redes só entra no último critério de desempate da cascata. No SNGS-113, as duas equipas jogam na mesma metade do campo, mas os dois guarda-redes distam mais de $70$ metros um do outro. Associar cada guarda-redes a um dos grupos permitiria determinar logo o lado correto.

    \item \textbf{Inclusão dos árbitros no agrupamento.} O agrupamento depende hoje do classificador de função para excluir a arbitragem. Em alternativa, poderia procurar cinco grupos, as duas equipas, os árbitros e os dois guarda-redes, à semelhança da solução vencedora de 2024 \citep{golovkin2025minimap}. Esta é a via para os $20$ vídeos da equipa de amarelo, que limitam a validação em ambos os sistemas.

    \item \textbf{Corte dos \emph{tracklets} nas trocas de jogador.} Alguns \emph{tracklets} seguem um jogador e, depois, outro, levando consigo a equipa ou o número do primeiro. Corrigir este problema do lado do número não melhorou a métrica, uma vez que a causa está no \emph{tracking}. A leitura por janela com confiança, que já existe, poderia indicar o ponto de troca, através de uma mudança no número lido com segurança.

    \item \textbf{Redução do custo computacional.} O sistema demora cerca de $22$ minutos por cada $30$ segundos de vídeo, sobretudo na reidentificação, na leitura do número e na estimativa dos pontos do campo. A confiança do modelo pode ser registada logo na primeira leitura, e o número pode ser lido apenas nos \emph{tracklets} que ainda não o têm. A reidentificação pode também ser limitada aos momentos em que o \emph{tracking} é ambíguo.

    \item \textbf{Vídeos de jogos não profissionais.} Os vídeos do SoccerNet-GSR vêm todos de transmissões televisivas de jogos profissionais. Em jogos amadores, a imagem pode ter menos resolução, os equipamentos podem ser menos distintos e o campo pode ter marcações incompletas, o que afeta a calibração. Ensaios preliminares com vídeos amadores, feitos fora desta dissertação, mostraram equipas mal separadas e árbitros dados como jogadores. Avaliar o sistema nestes vídeos, a começar por um conjunto anotado próprio, é o passo necessário para o usar fora do SoccerNet-GSR.
\end{enumerate}

Para além destas linhas, existe uma direção de outra natureza. Considerando apenas a posição, o sistema proposto não ultrapassa os $61.8\%$ de GS-HOTA no conjunto de teste. Os sistemas do topo do desafio ultrapassam os $60\%$ na métrica completa porque melhoram a deteção, o \emph{tracking} e a calibração, os módulos que este trabalho manteve intactos. Este é um âmbito distinto do que aqui foi definido.

Por fim, a direção seguida coincide com a adotada pelas melhores soluções da edição de 2025. A solução vencedora também recorre a um modelo de visão e linguagem para ler o número, o LLaMA-3.2-Vision, cuja versão mais pequena tem $11$ mil milhões de parâmetros. Neste trabalho, chegou-se à mesma ideia de forma independente, com um modelo de $3$ mil milhões de parâmetros que cabe numa placa gráfica de $6$~GB.
